\begin{abstract}
This article continues two analytic programs in the ProveIt polylogarithm
manuscript: sharp Euler remainder estimates at arbitrary real orders,
and the motion of real zeros under harmonic and Lerch deformation.
For the optimal normalized remainder $C_N$, we prove eventual exact
reduction to the outer-order boundary $a=0$, uniqueness of the optimizing
inner order, and the conjectured law
$C_N=1+cN^{-p}+O(N^{-p-\delta})$, with explicit positive constants.
For the pointwise remainder kernel, we prove global convergence of all
maximizers to a unique interior scaled limit, a convergent expansion in
$1/N$, and eventual strict decrease of its optimum. The first later
truncation, $N=2$, has an exact algebraic optimum of degree seven.
For the inverse equation $e^U=A+tU$, an exact rectangular barrier selects
the relevant analytic sheet and its two dominant square-root
singularities for every $A>1$. This makes the repository's conditional
velocity asymptotic unconditional, yields all-order oscillatory
expansions, bounded sign gaps, convergent boundary identities, and an
explicit limiting distribution of the zero-velocity lattice parameters.
The inverse series is identified with earlier work of Kalugin and
Jeffrey, and the coefficient polynomials with Cohen's Lambert
polynomials. For the latter we prove real-rootedness, strict interlacing,
an all-orders fixed-parameter expansion of the largest roots with the
polynomial coefficient degrees conjectured by Cohen, and a Bessel
limit at the smallest-root endpoint. The first two largest-root
corrections are explicit combinations of zeta values. All claims are
separated from numerical evidence and from
questions that remain open. Exact algebraic verification, numerical
diagnostics, source provenance, and integration notes accompany the
article.
\end{abstract}

\tableofcontents
\clearpage
\pagenumbering{arabic}

\section{Research scope, results, and relation to the baseline}
\label{sec:introduction}

The organizing question is how a global extremum or a global inverse
sheet can be recovered from a local asymptotic calculation. In both
problems considered here, the local calculation was substantially
available in the source material, but it did not by itself settle the
global conclusion. For Euler remainders, a trial family on the axis
gave the right candidate correction to the optimal constant; one still
had to exclude all other order parameters. For harmonic zero velocities,
formal critical points and local square-root expansions were known;
one still had to prove that the selected inverse branch actually reaches
exactly the proposed dominant pair. The proofs below supply those
global steps.

\subsection{A fixed and reproducible source comparison}

The repository baseline is revision
\texttt{4c173c06cc32c9cea554b39be1837ad2ae897fc1},
dated October 10, 2026. The review covered the modular manuscript,
its editorial ledger and validation notes, the harmonic-order geometry
continuation, and all six archives then present in \src{docs/incoming}.
The provenance manifest records the paths and Git blob hashes of the
retrieved sources. Statements about what was already proved or left
open refer to this revision, rather than to an evolving branch
\cite{ProveItSnapshot,HarmonicBaseline,UniformBoundsIncoming,SharpEulerIncoming}.

Several substantial problems have already been settled in that baseline.
In particular, the sharp constant uniform in \emph{all} positive real
orders and \emph{all} truncation indices is already known; so are the
$S_2$ and $S_4$ reductions and the certified radial double-extremum
phenomenon. They are not reported here as fresh results. The unresolved
$S_6$ reduction is not claimed. The selected targets are the remaining
global large-index Euler problem and the dominant-sheet question for
diagonal harmonic velocities, together with consequences naturally
forced by their proofs.

\subsection{The main analytic conclusions}

Use the precise normalization of \cref{sec:setup}. The first result is
\[
 C_N=\max_{b>0}R_N(0,b)\quad\hbox{for all sufficiently large }N,
\]
with a unique maximizing order $b_N$. Put
\[
 d=\log(3/2),\qquad p=\frac{\log2}{d},\qquad
 q=\frac{\log3}{\log2},\qquad
 c=\frac{p^p}{(p+1)^{p+1}},\qquad
 \delta=\frac{\log(10/9)}d.
\]
Then
\begin{align}
 b_N&=\frac{\log(Nq)}d+O(N^{-\delta}),\label{eq:intro-b}\\
 C_N&=1+cN^{-p}+O(N^{-p-\delta}).\label{eq:intro-C}
\end{align}
Here $p=1.70951129135\ldots$ and $c=0.167947796543\ldots$.
The positive-quadrant supremum is approached as $a\downarrow0$ and is
not attained there. The proof also gives a necessary scale for any
near-optimal outer order, sharpened to a necessary and sufficient
criterion and a joint profile with exact linear loss $\sqrt\pi$.
It uses probability kernels and positive
Taylor remainders, so no cancellation-sensitive finite Euler sum is
needed as a proof premise.

The pointwise kernel optimum $M_N$ is a different quantity. Its limit is
\[
 M_\infty=1.053861930362711626\ldots>1,
\]
whereas $C_N\to1$. Every global kernel maximizer satisfies
$Nr_N\to3.941064831647250668\ldots$ and
$y_N\to0.145303859259358424\ldots$. The optimized kernel value and
location have convergent expansions in $1/N$ near zero. The exact
degree-seven answer for $M_2$ gives a finite-index anchor for this
asymptotic theory.

For the harmonic problem, let $A=1+h>1$ encode the lattice spacing.
There is a unique $\theta\in(0,\pi)$ such that
\begin{equation}\label{eq:intro-angle}
 \log A=1-\theta\cot\theta+\log\frac\theta{\sin\theta}.
\end{equation}
Writing $a_*=1-\theta\cot\theta$, $u_*=a_*+i\theta$, and
$R=e^{a_*}$, the selected inverse of $e^U=A+tU$ has precisely the
two singularities $Re^{\pm i\theta}$ on its convergence circle.
If $v_n(A)$ is the initial velocity of the $n$th diagonal zero, then
\[
 v_n(A)=\frac{R^{-n}}{\sqrt\pi\,n^{3/2}}
          \Re\!\left(e^{-in\theta}\sqrt{-2u_*}\right)
          +O_A(R^{-n}n^{-5/2}),
\]
with an explicitly selected square root. The article gives its next
coefficient and an algorithm for every subsequent one. The error is
additive; no relative approximation is asserted near a cosine zero.

The same phase describes the distribution of the positive zeros of
the velocity polynomials. In the angle coordinate their limiting
distribution is uniform on $(0,\pi)$. In the logarithmic coordinate
$L=\log A$ its density is
\[
 \frac{\theta}{\pi\bigl((1-\theta\cot\theta)^2+\theta^2\bigr)}.
\]
Thus complex singularity geometry determines both oscillation in the
index and the distribution of real parameters where an initial
velocity vanishes.

The Lambert-polynomial identification also leads to an all-orders
solution of the fixed-parameter extreme-root questions in Cohen's
Conjecture 7.1 \cite{Cohen2021}. If $X_{n,k,j}$ is the $j$th largest
positive root of $L_{n,k}$, then for fixed positive integers $j,k$ and
any fixed $M$,
\begin{equation}\label{eq:intro-Cohen}
 X_{n,k,j}=\frac{n+k-1}{j}+\log n+\gamma-H_{j-1}
       +\sum_{m=1}^{M}\frac{C_{j,m}(k-1)}{n^m}
       +O_{j,k,M}(n^{-M-1}).
\end{equation}
Each $C_{j,m}$ is a polynomial of degree exactly $m-1$. The first
largest-root coefficient is
\[
 C_{1,1}=-\frac32-\zeta(2)-\zeta(3),
\]
and the next coefficient is given explicitly in
\cref{sec:cohen-edges}. Euler's constant occurs only in the leading
shift; every later coefficient lies in a specified algebra generated
by ordinary zeta values. A different endpoint scaling gives
$n^2x_{n,k,j}\to z_{k,j}^2/2$ for the $j$th smallest positive root,
where $z_{k,j}$ is a positive zero of $J_k$. The bulk, large-root,
and Bessel limits have separately stated regimes.

\subsection{Attribution and the meaning of a result here}

Kalugin and Jeffrey's 2012 Lambert-$W$ series already contains the
inverse coefficient array and the convergence-radius criterion after
an explicit change of variables \cite{KaluginJeffrey2012}. Cohen's
Lambert polynomials are the same coefficient polynomials up to a sign
\cite{Cohen2021}. These are substantive identifications, not merely
analogies. They correct a gap in the source comparison and make it
possible to connect the repository's questions to the existing
literature. The barrier proof, its detailed continuation statement,
and the derived velocity and zero-distribution results are presented
with this antecedent made explicit.

Throughout the article, ``proved'' means proved in ordinary mathematics
with the hypotheses stated. The exact polynomial computations have
independent finite verifiers. Numerical calculations illustrate scales,
locate examples, and test formulas; they are not substituted for
global analytic arguments. No assertion of literature-wide priority
or machine-checked formalization is made. The final sections give
precise proposed source amendments and distinguish completed targets
from the remaining finite-index, edge-scaling, and arithmetic questions.

